% ============================================================
% PARTIE I - CHAPITRE I : La structure d'accueil et le sujet
% ============================================================
\chapter{La structure d'accueil et le sujet}
\pagestyle{styleGlobal}
\label{chap:structure_sujet}

% ============================================================
\section{La structure d'accueil}
% ============================================================

\subsection{Historique}

Akieni est une société spécialisée dans la transformation numérique et les solutions logicielles en République du Congo. Implantée à Brazzaville depuis 2022, elle a pour mission d'accompagner les entreprises et les institutions publiques et privées dans leurs projets de transformation digitale. Elle est reconnue comme un acteur majeur du numérique et collabore avec des gouvernements, des institutions, des entreprises et des citoyens à travers l’Afrique centrale. 

Parmi ses initiatives, Akieni a lancé le programme de formation \textit{Akieni Academy}, dont la première promotion a démarré officiellement le 15 novembre 2024 à Brazzaville. À cette occasion, 120 apprenants ont été sélectionnés parmi 1 740 candidats présélectionnés en ligne, venant du Congo, du Cameroun, du Rwanda et de la République Démocratique du Congo. La deuxième promotion a été lancée en janvier 2026, avec l'ambition de former 400 nouveaux apprenants. L'objectif à long terme d'Akieni est de former 30 000 jeunes aux métiers du numérique d'ici 2030, afin de positionner le Congo comme un pôle d'exportation de compétences digitales. La structure est dirigée par M. Frédéric NZÉ, actuel Ministre des Postes, des Télécommunications et de l'Économie numérique de la République du Congo.

\subsection{Missions}

Akieni a plusieurs missions principales :

\begin{itemize}
	\item accompagner les organisations publiques et privées dans leurs projets de transformation digitale ;
	\item promouvoir les talents locaux et organiser des événements numériques tels que des hackathons pour encourager l'innovation ;
	
	\item investir dans le capital humain ;
	\item stimuler la transformation par la collaboration ;
	\item favoriser l'innovation numérique.
\end{itemize}

\subsection{Organigramme Général}

La figure~\ref{fig:organigramme_akieni} présente l'organisation générale d'Akieni, structurée autour d'une Direction Générale qui supervise cinq directions principales, chacune pilotant ses propres départements opérationnels.

\begin{figure}[H]
	\centering
	\includegraphics[width=0.95\textwidth]{organigramme_akieni.png}
	\caption{Organigramme général d'Akieni}
	\label{fig:organigramme_akieni}
\end{figure}

\subsection{Attributions des structures}

Akieni est organisée en cinq directions principales, chacune ayant un domaine de responsabilité bien défini.

\textbf{La Direction des Ressources Humaines} gère le recrutement, la gestion des carrières et le bien-être du personnel. Elle veille à ce que les équipes disposent des compétences nécessaires pour mener à bien les projets de la structure.

\textbf{La Direction Technique} est le cœur opérationnel d'Akieni sur le plan technologique. Elle regroupe trois départements : le département de développement Front-end et Back-end, qui conçoit et maintient les applications et plateformes numériques ; le département DevOps, qui assure l'intégration continue, le déploiement et la supervision des infrastructures ; et le département Support, qui gère l'assistance technique auprès des utilisateurs internes et des clients.

\textbf{La Direction Administrative et Financière} assure la gestion des ressources financières et administratives de la structure. Elle comprend le département Finance, chargé de la comptabilité et du suivi budgétaire ; le département Moyens Généraux, qui gère les équipements et les locaux ; et le département Contrôle de Gestion, qui évalue la performance financière de l'organisation.

\textbf{La Direction Programme et Solution} pilote les projets clients et les offres de services d'Akieni. Elle s'appuie sur quatre départements : le PMO (\textit{Project Management Office}), qui coordonne et standardise la gestion des projets ; le département Presale, qui accompagne les avant-ventes et la rédaction des offres commerciales ; le département Consulting, qui propose des conseils en transformation numérique ; et le département Customer Excellence, qui veille à la satisfaction client et à l'excellence de la qualité des livrables.

\textbf{La Direction de la Communication} est responsable de la communication interne et externe de la structure, de la gestion de l'image de marque d'Akieni, ainsi que de la promotion de ses programmes et de ses offres auprès du grand public et des partenaires.

\subsection{Situation informatique}

\subsubsection{Personnel informatique}

Akieni étant une structure spécialisée dans le numérique, l'essentiel de son personnel est à profil technique. L'équipe informatique est rattachée à la Direction Technique et se compose notamment de développeurs Full Stack (Front-end et Back-end), d'ingénieurs DevOps chargés de l'infrastructure et du déploiement, et de techniciens support. La Direction Programme et Solution compte également des consultants et des chefs de projet qui encadrent les missions techniques chez les clients.

\subsubsection{Matériels informatiques}

Pour mener ses activités de formation, de développement et de conseil, Akieni dispose d'un parc matériel adapté à ses besoins. On y trouve notamment des postes de travail et des ordinateurs portables mis à disposition des équipes techniques et des apprenants, des serveurs locaux utilisés pour les environnements de développement et de test, ainsi que les équipements réseau nécessaires à la connectivité de ses locaux (routeurs, switchs, points d'accès Wi-Fi). Les projections et présentations en salle de formation sont assurées par des vidéoprojecteurs et des écrans interactifs.

\subsubsection{Logiciels informatiques}

Akieni utilise un ensemble d'outils logiciels couvrant ses différents besoins. Sur le plan du développement, les équipes travaillent avec des environnements tels que Visual Studio Code, Git et GitHub pour la gestion du code source, ainsi que des plateformes cloud comme AWS ou Azure pour l'hébergement des applications. Pour la gestion de projets, des outils comme Jira ou Trello sont utilisés au sein du PMO. La communication interne repose sur des solutions telles que Slack ou Microsoft Teams. Enfin, les formations de l'Akieni Academy s'appuient sur des plateformes d'apprentissage en ligne (LMS) permettant de suivre la progression des apprenants à distance.

% ============================================================
\section{Le sujet}
% ============================================================

\subsection{Contexte du sujet}

Le numérique occupe aujourd'hui une place centrale dans l'organisation des services professionnels. Les plateformes numériques ont remplacé progressivement les méthodes traditionnelles de recherche de prestataires, offrant plus de transparence et d'efficacité. Cependant, à Brazzaville, de nombreux professionnels exercent dans des domaines variés comme l'artisanat, les services techniques, les prestations intellectuelles ou les services à la personne, mais ils restent peu visibles faute d'une plateforme numérique adaptée.

De leur côté, les particuliers, les entreprises et les institutions ont du mal à trouver rapidement un prestataire compétent, disponible et de confiance. Les échanges reposent encore largement sur le bouche-à-oreille, sans aucune garantie de qualité ni de sécurité financière.

« SALISA » naît de ce constat. Il s'agit d'une application cross-plateforme conçue pour faciliter la mise en relation entre prestataires et clients, en intégrant un système de vérification des profils, un mécanisme de paiement sécurisé et un outil de notation des prestations. Le nom \textit{SALISA}, emprunté au lingala, a une double signification : pour le prestataire, \textit{ko salisa} signifie \textit{"aider"} et pour le demandeur, \textit{salisa} signifie \textit{"faire faire"} (un travail).

\subsection{Problématique du sujet}

Malgré la présence de nombreux prestataires de services à Brazzaville, il n'existe pas de mécanisme fiable permettant de les mettre en relation avec les clients. Les échanges se font de façon informelle, sans garantie de qualité, de ponctualité ni de sécurité financière pour les parties concernées.

La problématique principale de ce travail peut être formulée ainsi :

\begin{center}
	\itshape
	Comment concevoir et mettre en \oe{}uvre une application cross-plateforme efficace permettant de connecter des prestataires et des clients, tout en garantissant la fiabilité des services, la sécurité des transactions et une expérience utilisateur adaptée au contexte local ?
\end{center}

Cette problématique soulève plusieurs enjeux : la gestion de la confiance entre utilisateurs, l'optimisation du choix des prestataires grâce à un mécanisme de mise en correspondance intelligent, la sécurisation des paiements, et l'accessibilité de la solution sur des appareils variés.

\subsection{Intérêts du sujet}

L'intérêt de ce travail se situe à plusieurs niveaux.

\begin{itemize}
	\item[\textbf{-}] \textbf{Sur le plan économique :} la plateforme permet de dynamiser le marché des services en facilitant l'accès aux opportunités pour les prestataires et en réduisant les difficultés de recherche pour les clients. Elle contribue à la création de valeur économique locale et à l'insertion professionnelle.
	
	\item[\textbf{-}] \textbf{Sur le plan social :} « SALISA » favorise l'inclusion numérique en offrant à un large public la possibilité d'accéder à des services fiables et de faire connaître leurs compétences. Elle encourage également les transactions formelles entre citoyens.
	
	\item[\textbf{-}] \textbf{Sur le plan technologique :} le travail intègre des solutions modernes comme la mise en correspondance intelligente (matching), la géolocalisation par quartier et le paiement mobile via MTN MoMo et Airtel Money, ce qui renforce sa pertinence dans un contexte numérique en développement.
	
	\item[\textbf{-}] \textbf{Sur le plan académique :} ce travail est une application concrète des enseignements du génie logiciel, notamment l'analyse des besoins, la modélisation UML, la conception de systèmes distribués et le développement d'applications selon la démarche \textbf{2TUP}.
\end{itemize}

% ============================================================
\section{Concepts liés au sujet}

La compréhension de ce travail nécessite la maîtrise de plusieurs concepts clés, définis ci-après.

\begin{itemize}
	\item[\textbf{-}] \textbf{Application cross-plateforme :} désigne une application conçue pour fonctionner sur plusieurs types de terminaux et de systèmes d'exploitation (web, Android, iOS) à partir d'une seule base de code, ce qui réduit les coûts et les délais de développement.
	
	\item[\textbf{-}] \textbf{Mise en relation :} processus qui consiste à connecter une offre de service proposée par un prestataire à une demande exprimée par un client, en s'appuyant sur des critères de compatibilité définis.
	
	\item[\textbf{-}] \textbf{Matching :} mécanisme qui identifie automatiquement les prestataires les plus adaptés à une demande, en tenant compte des compétences, de la localisation géographique, de la disponibilité, du budget et des évaluations reçues.
	
	\item[\textbf{-}] \textbf{Système de confiance :} ensemble de mécanismes tels que les notes, les avis, la vérification d'identité et les badges de certification, qui permettent de s'assurer de la fiabilité des utilisateurs et de la qualité des services proposés.
	
	\item[\textbf{-}] \textbf{Paiement mobile :} solution de paiement électronique très utilisée en Afrique, qui permet d'effectuer des transactions financières via un téléphone mobile, sans avoir besoin d'un compte bancaire. Au Congo, les deux opérateurs de paiement mobile sont MTN MoMo et Airtel Money.
	
	\item[\textbf{-}] \textbf{API (Application Programming Interface) :} interface de programmation qui permet à deux applications de communiquer entre elles en échangeant des données selon un protocole défini. Dans le cadre de « SALISA », les API de MTN MoMo et d'Airtel Money sont utilisées pour intégrer le paiement mobile directement dans la plateforme.
	
	\item[\textbf{-}] \textbf{Progressive Web Application (PWA) :} application web qui offre une expérience proche d'une application mobile classique, avec un fonctionnement possible hors connexion, l'installation sur l'écran d'accueil et des notifications. Elle ne nécessite pas de téléchargement depuis un store d'applications.
	
	\item[\textbf{-}] \textbf{UML (Unified Modeling Language) :} langage de modélisation graphique utilisé en génie logiciel pour représenter les différents aspects d'un système informatique, comme les cas d'utilisation, les classes, les séquences ou le déploiement.
	
	\item[\textbf{-}] \textbf{2TUP (Two Tracks Unified Process) :} processus de développement logiciel dérivé du Processus Unifié (UP), qui organise la conception selon deux axes parallèles : la branche fonctionnelle (besoins métier) et la branche technique (contraintes d'architecture), qui fusionnent en phase de réalisation.
	
	\item[\textbf{-}] \textbf{Géolocalisation :} technique qui permet de situer un utilisateur ou un service dans un espace géographique, afin de faciliter la mise en relation de proximité entre prestataires et clients.
\end{itemize}